\documentclass[11pt]{article}

\usepackage[margin=1in]{geometry}
\usepackage{amsmath,amssymb,amsthm}
\usepackage{booktabs}
\usepackage{array}
\usepackage{hyperref}

\newtheorem{theorem}{Theorem}
\newtheorem{proposition}{Proposition}
\newtheorem{remark}{Remark}

\newcommand{\R}{\mathbb{R}}
\newcommand{\E}{\mathbb{E}}
\newcommand{\Prb}{\mathbb{P}}
\newcommand{\HK}{\mathrm{HK}}
\newcommand{\OPT}{\mathrm{OPT}}
\newcommand{\Agg}{\mathrm{Agg}}

\title{Scenario-Based and Aggregate LPs for A Posteriori TSP}
\author{}
\date{}

\begin{document}
\maketitle

\section{The a posteriori TSP model}

Let $r$ be a depot that is always present, and let $V$ be a set of stochastic vertices. Each vertex $v\in V$ is activated independently with probability $p_v$. For an activation set $A\subseteq V$, the realized vertex set is
\[
R_A := A\cup\{r\}.
\]
We assume that the distances are given by a symmetric metric $d$ on $V\cup\{r\}$. Thus
\[
d(u,v)=d(v,u), \qquad d(u,u)=0,
\]
and the triangle inequality holds. For a set $W\subseteq V\cup\{r\}$, let $E(W)$ denote the edge set of the complete graph on $W$. For $S\subseteq W$, let $\delta_W(S)$ denote the set of edges of $E(W)$ with exactly one endpoint in $S$. When the underlying vertex set is clear, we simply write $\delta(S)$.

For each $A\subseteq V$, define
\[
q_A := \Prb[\text{the active set is }A]
= \prod_{v\in A}p_v \prod_{v\in V\setminus A}(1-p_v).
\]

\section{The scenario-wise Held--Karp relaxation}

For a fixed active set $A\subseteq V$, the Held--Karp relaxation for the TSP instance on $R_A=A\cup\{r\}$ is
\[
\begin{aligned}
\HK(A):=\min \quad
& \sum_{e\in E(R_A)} d_e x^A_e \\
\text{s.t.}\quad
& x^A(\delta_{R_A}(v)) = 2
&& \forall v\in R_A, \\
& x^A(\delta_{R_A}(S)) \ge 2
&& \forall \emptyset\neq S\subsetneq R_A, \\
& x^A_e\ge 0
&& \forall e\in E(R_A).
\end{aligned}
\]
Here $d_e=d(u,v)$ for $e=uv$.

The expected scenario-wise Held--Karp value is
\[
\E_A[\HK(A)] = \sum_{A\subseteq V} q_A\HK(A).
\]

\section{An LP for the average of the scenario-wise Held--Karp solutions}

The following extended LP computes the average of the scenario-wise Held--Karp relaxations. It has a separate Held--Karp vector $x^A$ for every activation set $A\subseteq V$.
\[
\begin{aligned}
\min \quad
& \sum_{A\subseteq V} q_A \sum_{e\in E(R_A)} d_e x^A_e \\
\text{s.t.}\quad
& x^A(\delta_{R_A}(v)) = 2
&& \forall A\subseteq V,\ \forall v\in R_A, \\
& x^A(\delta_{R_A}(S)) \ge 2
&& \forall A\subseteq V,\ \forall \emptyset\neq S\subsetneq R_A, \\
& x^A_e\ge 0
&& \forall A\subseteq V,\ \forall e\in E(R_A).
\end{aligned}
\tag{SCEN}
\]
This LP is just the probability-weighted sum of the Held--Karp LPs for the individual scenarios. Therefore its optimum value is exactly
\[
\sum_{A\subseteq V}q_A\HK(A).
\]

If one wants to explicitly record the average edge-vector, introduce variables $y_e$ for $e\in E(V\cup\{r\})$ and add the equations
\[
y_e=\sum_{A:e\in E(R_A)}q_A x^A_e
\qquad \forall e\in E(V\cup\{r\}).
\]
The vector $y$ is then the expected Held--Karp solution:
\[
y=\E_A[x^A].
\]

\section{The aggregate cut LP}

A natural attempt is to describe the average vector $y$ directly using only averaged degree and cut constraints. This gives the following aggregate LP.

For $\emptyset\neq S\subsetneq V\cup\{r\}$, define
\[
T(S):=
\begin{cases}
S, & r\notin S,\\
(V\cup\{r\})\setminus S, & r\in S.
\end{cases}
\]
Then $T(S)\subseteq V$, and the cut $S$ is crossed by the realized vertex set $R_A=A\cup\{r\}$ if and only if $A\cap T(S)\neq\emptyset$. Hence define
\[
\rho(S):=2\Prb[A\cap T(S)\neq\emptyset]
=2\left(1-\prod_{v\in T(S)}(1-p_v)\right).
\]

The aggregate LP is
\[
\begin{aligned}
\min \quad
& \sum_{e\in E(V\cup\{r\})} d_e y_e \\
\text{s.t.}\quad
& y(\delta(v))=2p_v
&& \forall v\in V, \\
& y(\delta(r))=2\Prb[A\neq\emptyset]
=2\left(1-\prod_{v\in V}(1-p_v)\right), \\
& y(\delta(S))\ge \rho(S)
&& \forall \emptyset\neq S\subsetneq V\cup\{r\}, \\
& y_e\ge 0
&& \forall e\in E(V\cup\{r\}).
\end{aligned}
\tag{AGG}
\]

Every average of scenario-wise Held--Karp solutions is feasible for \((\mathrm{AGG})\). Indeed, if
\[
y_e=\sum_{A:e\in E(R_A)}q_Ax^A_e,
\]
where $x^A$ is feasible for $\HK(A)$, then the degree equations and the aggregate cut inequalities follow by taking expectations of the corresponding scenario-wise Held--Karp constraints.

However, \((\mathrm{AGG})\) is not an exact description of the projection of \((\mathrm{SCEN})\). It only keeps the averaged cut inequalities and forgets the requirement that the vector $y$ must decompose into scenario-wise feasible Held--Karp vectors.

\section{A counterexample}

We now give an explicit counterexample showing that \((\mathrm{AGG})\) can be strictly weaker than the scenario-indexed LP \((\mathrm{SCEN})\).

Let the depot be $0$, and let the stochastic vertex set be
\[
V=\{1,2,3,4\}.
\]
Each stochastic vertex is independently active with probability $1/2$. Hence every scenario has probability $1/16$.

Use the following symmetric integral metric $d$ on $\{0,1,2,3,4\}$:
\[
\begin{array}{c|ccccc}
 & 0 & 1 & 2 & 3 & 4 \\
\hline
0 & 0 & 14 & 13 & 8 & 7 \\
1 & 14 & 0 & 16 & 6 & 10 \\
2 & 13 & 16 & 0 & 10 & 12 \\
3 & 8 & 6 & 10 & 0 & 11 \\
4 & 7 & 10 & 12 & 11 & 0
\end{array}
\]

\subsection{The aggregate LP solution}

The following vector is feasible for \((\mathrm{AGG})\):
\[
\begin{array}{c|cccccccccc}
e & 01&02&03&04&12&13&14&23&24&34\\
\hline
y_e
& \frac{5}{16}
& \frac{7}{16}
& \frac{9}{16}
& \frac{9}{16}
& \frac{3}{16}
& \frac{4}{16}
& \frac{4}{16}
& \frac{3}{16}
& \frac{3}{16}
& \frac{0}{16}
\end{array}
\]
The degree constraints for the stochastic vertices are all tight:
\[
y(\delta(v))=1=2\cdot \frac12
\qquad \forall v\in\{1,2,3,4\}.
\]
The depot degree constraint is also tight:
\[
y(\delta(0))
=\frac{5+7+9+9}{16}
=\frac{30}{16}
=\frac{15}{8}
=2\left(1-\frac1{16}\right).
\]
The aggregate cut inequalities are also satisfied. Its objective value is
\[
\begin{aligned}
d\cdot y
&=14\cdot\frac{5}{16}
+13\cdot\frac{7}{16}
+8\cdot\frac{9}{16}
+7\cdot\frac{9}{16} \\
&\qquad
+16\cdot\frac{3}{16}
+6\cdot\frac{4}{16}
+10\cdot\frac{4}{16}
+10\cdot\frac{3}{16}
+12\cdot\frac{3}{16}
+11\cdot\frac{0}{16} \\
&=29.625
=\frac{237}{8}.
\end{aligned}
\]
In fact this vector is optimal for \((\mathrm{AGG})\), so
\[
\OPT_{\Agg}=\frac{237}{8}.
\]

\subsection{The scenario-wise Held--Karp values}

The scenario-wise Held--Karp values are as follows.
\[
\begin{array}{c|c|c}
A & \HK(A) & \text{one optimal tour} \\
\hline
\emptyset & 0 & (0) \\
\{1\} & 28 & 0-1-0 \\
\{2\} & 26 & 0-2-0 \\
\{3\} & 16 & 0-3-0 \\
\{4\} & 14 & 0-4-0 \\
\{1,2\} & 43 & 0-1-2-0 \\
\{1,3\} & 28 & 0-1-3-0 \\
\{1,4\} & 31 & 0-1-4-0 \\
\{2,3\} & 31 & 0-2-3-0 \\
\{2,4\} & 32 & 0-2-4-0 \\
\{3,4\} & 26 & 0-3-4-0 \\
\{1,2,3\} & 43 & 0-1-3-2-0 \\
\{1,2,4\} & 46 & 0-2-1-4-0 \\
\{1,3,4\} & 31 & 0-3-1-4-0 \\
\{2,3,4\} & 37 & 0-3-2-4-0 \\
\{1,2,3,4\} & 46 & 0-2-3-1-4-0
\end{array}
\]
Therefore
\[
\sum_{A\subseteq V}q_A\HK(A)
=\frac1{16}\sum_{A\subseteq V}\HK(A)
=29.875
=\frac{239}{8}.
\]
Thus the scenario-indexed LP has value
\[
\OPT_{\mathrm{SCEN}}=\frac{239}{8}.
\]

On the other hand, the aggregate LP has value
\[
\OPT_{\Agg}=\frac{237}{8}.
\]
Consequently,
\[
\OPT_{\Agg}<\OPT_{\mathrm{SCEN}}.
\]
This proves that the aggregate cut LP is not an exact description of the average of the scenario-based Held--Karp LPs.

\begin{remark}
The reason for the failure is structural. The aggregate LP retains only inequalities obtained by averaging individual Held--Karp cut constraints. It does not enforce that the aggregate vector $y$ decomposes as
\[
y=\sum_{A\subseteq V}q_Ax^A
\]
with each $x^A$ feasible for the Held--Karp relaxation on $A\cup\{r\}$. Thus the aggregate LP can combine cut-capacity mass in a way that is feasible after averaging but cannot be realized scenario by scenario.
\end{remark}

\end{document}
